% ====================================================================
% ANALISE-REQUISITOS.TEX — Seção 2: Análise de Requisitos
% ====================================================================

\section{Análise de Requisitos}

Esta seção apresenta os requisitos levantados a partir da descrição do projeto fornecida
pelo cliente, organizados em requisitos funcionais e não funcionais.

% --------------------------------------------------------------------
\subsection{Requisitos Funcionais}
% --------------------------------------------------------------------
% Requisitos funcionais descrevem O QUE o sistema deve fazer.

Os requisitos funcionais definem os comportamentos e funcionalidades que o sistema
deve apresentar.

\subsubsection{Requisitos de Prioridade 1 (Entrega na Etapa de Protótipo)}

\begin{itemize}
    \item \textbf{RF01 — Leitura de temperatura e umidade:} O sistema deve coletar
    dados de temperatura e umidade do ambiente por meio de um sensor dedicado.

    \item \textbf{RF02 — Exibição em tempo real:} Os dados coletados devem ser
    disponibilizados em tempo real na interface web.

    \item \textbf{RF03 — Modo Desligado:} O sistema deve suportar um modo de operação
    no qual o sensor e o sistema são completamente desativados.

    \item \textbf{RF04 — Modo Ligado:} O sistema deve suportar um modo de operação
    normal, aguardando ação do usuário para realizar leituras.

    \item \textbf{RF05 — Modo Automático:} O sistema deve suportar um modo de operação
    no qual a coleta e exibição dos dados ocorrem de forma contínua e automática, sem
    intervenção do usuário.

    \item \textbf{RF06 — Alternância de modos pela interface web:} A interface web deve
    disponibilizar botões para que o usuário alterne entre os três modos de operação
    (Desligado, Ligado e Automático).

    \item \textbf{RF07 — Indicação visual por LEDs:} O hardware deve contar com LEDs
    indicadores que reflitam o modo de operação ativo, acendendo ou apagando conforme
    o estado atual do sistema.
\end{itemize}

\subsubsection{Requisitos de Prioridade 2 (Extras Opcionais)}

\begin{itemize}
    \item \textbf{RF08 — Histórico de leituras:} A interface web pode exibir um
    histórico de leituras com data e hora das medições.

    \item \textbf{RF09 — Alertas automáticos:} O sistema pode emitir alertas
    automáticos na interface web quando os valores de temperatura ou umidade
    ultrapassarem limites pré-definidos.

    \item \textbf{RF10 — Notificação por e-mail:} O sistema pode enviar notificações
    por e-mail ao usuário quando os valores monitorados ultrapassarem os limites
    configurados.

    \item \textbf{RF11 — Log local:} O sistema pode armazenar os logs de leituras em
    arquivo \texttt{.txt}, garantindo persistência de dados localmente.

    \item \textbf{RF12 — Aplicativo mobile:} O sistema pode disponibilizar um aplicativo
    mobile (Android e/ou iOS) como alternativa à interface web para monitoramento e
    controle do sistema.
\end{itemize}

% --------------------------------------------------------------------
\subsection{Requisitos Não Funcionais}
% --------------------------------------------------------------------
% Requisitos não funcionais descrevem COMO o sistema deve se comportar
% (desempenho, confiabilidade, usabilidade, restrições técnicas etc.)

Os requisitos não funcionais definem as qualidades e restrições técnicas do sistema.

\begin{itemize}
    \item \textbf{RNF01 — Comunicação em rede:} O microcontrolador deve ser capaz
    de se comunicar com a interface web via rede (Wi-Fi ou equivalente).

    \item \textbf{RNF02 — Justificativa técnica de componentes:} A seleção do sensor
    e do microcontrolador deve ser justificada com base em características técnicas
    relevantes, tais como faixa de medição, precisão, consumo de energia, facilidade
    de programação e suporte a comunicação via rede.

    \item \textbf{RNF03 — Atualização em tempo real:} A interface web deve exibir
    os dados de temperatura e umidade sem necessidade de recarregar a página
    (atualização dinâmica).

    \item \textbf{RNF04 — Usabilidade:} A interface web deve ser simples e intuitiva,
    permitindo ao usuário alternar entre os modos de operação sem dificuldade.
\end{itemize}
